\begin{figure}[htbp]
\centering
\resizebox{.4836\textwidth}{!}{%
\begin{circuitikz}
\tikzstyle{every node}=[font=\LARGE]

%% ===================================================================
%% HAUPTKAMMER - Außenwände und Grundstruktur
%% ===================================================================
\draw [line width=1.5pt, short] (1.25,16) -- (2.25,16);
\draw [line width=1.5pt, short] (2.25,16) -- (2.25,7);
\draw [line width=1.5pt, short] (20.25,7) -- (20.25,16);
\draw [line width=1.5pt, short] (20.25,16) -- (21.25,16);
\draw [line width=1.5pt, short] (1.25,16) -- (1.25,16.75);
\draw [line width=1.5pt, short] (1.25,16.75) -- (1.5,16.75);
\draw [line width=1.5pt, short] (2.25,16.75) -- (2.5,16.75);
\draw [line width=1.5pt, short] (2.5,16.75) -- (2.5,7);
\draw [line width=1.5pt, short] (2.5,7) -- (20,7);
\draw [line width=1.5pt, short] (20,7) -- (20,16.75);
\draw [line width=1.5pt, short] (21.25,16) -- (21.25,16.75);
\draw [line width=1.5pt, short] (21.25,16.75) -- (21,16.75);
\draw [line width=1.5pt, short] (20,16.75) -- (20.25,16.75);

%% ===================================================================
%% ANSCHLÜSSE OBEN - Verschraubungen und Verbindungsstücke
%% ===================================================================
\draw [ line width=1.5pt ] (1.75,16.75) circle (0.25cm);
\draw [ line width=1.5pt ] (20.5,16.75) circle (0.25cm);
\draw [line width=1.5pt, short] (2.25,16.75) -- (2,16.75);
\draw [line width=1.5pt, short] (21,16.75) -- (20.75,16.75);
\draw [line width=1.5pt, short] (1.25,17) -- (2.75,17);
\draw [line width=1.5pt, short] (2.75,17) -- (2.75,16.25);
\draw [line width=1.5pt, short] (21.25,18.25) -- (21.25,18);
\draw [line width=1.5pt, short] (2,16.75) -- (2,16.5);
\draw [line width=1.5pt, short] (2,16.5) -- (1.5,16.5);
\draw [line width=1.5pt, short] (1.5,16.75) -- (1.5,16.5);
\draw [line width=1.5pt, short] (20.75,16.75) -- (20.75,16.5);
\draw [line width=1.5pt, short] (20.75,16.5) -- (20.25,16.5);
\draw [line width=1.5pt, short] (20.25,16.75) -- (20.25,16.5);
\draw [line width=1.5pt, short] (1.5,17) -- (1.5,17);
\draw [line width=1.5pt, short] (1.5,17) -- (1.5,17);
\draw [line width=1.5pt, short] (17.5,17) -- (17.5,17.75);
\draw [line width=1.5pt, short] (21.25,18) -- (17.75,18);
\draw [line width=1.5pt, short] (17.75,18) -- (17.5,17.75);
\draw [line width=1.5pt, short] (21.25,17) -- (20,17);
\draw [line width=1.5pt, short] (20,16.25) -- (19,16.25);
\draw [line width=1.5pt, short] (19,16.25) -- (19,16.75);
\draw [line width=1.5pt, short] (19,16.75) -- (18.75,17);
\draw [line width=1.5pt, short] (20,16.75) -- (20,17);
\draw [line width=1.5pt, short] (18.75,17) -- (18.75,17.25);
\draw [line width=1.5pt, short] (18.75,17.25) -- (21.25,17.25);
\draw [line width=1.5pt, short] (21.25,17.25) -- (21.25,17);
\draw [line width=1.5pt, short] (2.75,16.25) -- (3.5,16.25);
\draw [line width=1.5pt, short] (3.5,16.25) -- (3.5,16.75);
\draw [line width=1.5pt, short] (3.5,16.75) -- (3.75,17);

%% ===================================================================
%% KAMMERWÄNDE - Graue Füllungen (Wanddicke und Struktur)
%% ===================================================================
\draw [ color={rgb,255:red,192; green,192; blue,192} , fill={rgb,255:red,192; green,192; blue,192}, line width=0.2pt ] (20,16.75) rectangle (20.25,7);
\draw [ color={rgb,255:red,192; green,192; blue,192} , fill={rgb,255:red,192; green,192; blue,192}, line width=0.2pt ] (20,7) rectangle (20,6.75);
\draw [ color={rgb,255:red,192; green,192; blue,192} , fill={rgb,255:red,192; green,192; blue,192}, line width=0.2pt ] (1.25,16.75) rectangle (1.5,16);
\draw [ color={rgb,255:red,192; green,192; blue,192} , fill={rgb,255:red,192; green,192; blue,192}, line width=0.2pt ] (1.5,16.5) rectangle (2.25,16);
\draw [ color={rgb,255:red,192; green,192; blue,192} , fill={rgb,255:red,192; green,192; blue,192}, line width=0.2pt ] (2.25,16.75) rectangle (2.5,7);
\draw [ color={rgb,255:red,192; green,192; blue,192} , fill={rgb,255:red,192; green,192; blue,192}, line width=0.2pt ] (2,16.75) rectangle (2.25,16.5);
\draw [ color={rgb,255:red,192; green,192; blue,192} , fill={rgb,255:red,192; green,192; blue,192}, line width=0.2pt ] (20.75,16.75) rectangle (21.25,16);
\draw [ color={rgb,255:red,192; green,192; blue,192} , fill={rgb,255:red,192; green,192; blue,192}, line width=0.2pt ] (20.25,16.5) rectangle (20.75,16);
\draw [ color={rgb,255:red,192; green,192; blue,192} , fill={rgb,255:red,192; green,192; blue,192}, line width=0.2pt ] (20.25,17) rectangle (20,16.75);
\draw [ color={rgb,255:red,192; green,192; blue,192} , fill={rgb,255:red,192; green,192; blue,192}, line width=0.2pt ] (19,17) rectangle (20,16.25);
\draw [ color={rgb,255:red,192; green,192; blue,192} , fill={rgb,255:red,192; green,192; blue,192}, line width=0.2pt ] (21,17) rectangle (19,17.25);
\draw [short] (20,16.25) -- (20,17);
\draw [short] (20.5,17) -- (20,17);
\draw [ color={rgb,255:red,192; green,192; blue,192} , fill={rgb,255:red,192; green,192; blue,192}] (21,17.25) rectangle (21.25,17);
\draw [ color={rgb,255:red,192; green,192; blue,192} , fill={rgb,255:red,192; green,192; blue,192}] (18.75,17.25) rectangle (19,17);
\draw [ color={rgb,255:red,192; green,192; blue,192} , fill={rgb,255:red,192; green,192; blue,192}] (19.75,6.25) -- (19.5,6.25) -- (20,7) -- (20.25,7) -- cycle;

%% ===================================================================
%% BODEN - Untere Kammerbasis und Verbindungen
%% ===================================================================
\draw [ color={rgb,255:red,192; green,192; blue,192} , fill={rgb,255:red,192; green,192; blue,192}] (2.5,7) -- (2.25,7) -- (2.75,6.25) -- (3,6.25) -- cycle;
\draw [ color={rgb,255:red,192; green,192; blue,192} , fill={rgb,255:red,192; green,192; blue,192}] (2.5,7) rectangle (3,6.75);
\draw [ color={rgb,255:red,192; green,192; blue,192} , fill={rgb,255:red,192; green,192; blue,192}] (2.75,6.75) rectangle (3.25,6.5);
\draw [ color={rgb,255:red,192; green,192; blue,192} , fill={rgb,255:red,192; green,192; blue,192}] (2.75,6.75) rectangle (3,6.5);
\draw [ color={rgb,255:red,192; green,192; blue,192} , fill={rgb,255:red,192; green,192; blue,192}] (3,6.75) rectangle (2.75,6.5);
\draw [ color={rgb,255:red,192; green,192; blue,192} , fill={rgb,255:red,192; green,192; blue,192}] (2.75,6.5) rectangle (3.25,6.25);
\draw [ color={rgb,255:red,192; green,192; blue,192} , fill={rgb,255:red,192; green,192; blue,192}] (2.75,6.75) rectangle (2.75,6.5);
\draw [ color={rgb,255:red,192; green,192; blue,192} , fill={rgb,255:red,192; green,192; blue,192}] (2.75,6.75) circle (0.25cm);
\draw [ color={rgb,255:red,192; green,192; blue,192} , fill={rgb,255:red,192; green,192; blue,192}] (3,7) rectangle (19.75,6.25);
\draw [ color={rgb,255:red,192; green,192; blue,192} , fill={rgb,255:red,192; green,192; blue,192}, line width=0.2pt ] (19.75,6.75) circle (0.25cm);
\draw [ color={rgb,255:red,192; green,192; blue,192} , fill={rgb,255:red,192; green,192; blue,192}, line width=0.2pt ] (19.75,7) rectangle (20,6.75);
\draw [line width=0.6pt, short] (2.75,6.25) -- (19.75,6.25);
\draw [line width=0.6pt, short] (19.75,6.25) -- (20.25,7);
\draw [line width=0.6pt, short] (2.25,7) -- (2.75,6.25);

%% ===================================================================
%% AUSGANGSVENTIL-BEREICH (rechts)
%% ===================================================================
\draw [ color={rgb,255:red,192; green,192; blue,192} , fill={rgb,255:red,192; green,192; blue,192}, line width=0.6pt ] (18.75,17) -- (19.25,17) -- (19.75,16.5) -- (19.25,16.5) -- cycle;
\draw [ color={rgb,255:red,192; green,192; blue,192} , fill={rgb,255:red,192; green,192; blue,192}, line width=0.6pt ] (19,17.25) rectangle (19.75,16.25);
\draw [ color={rgb,255:red,192; green,192; blue,192} , fill={rgb,255:red,192; green,192; blue,192}, line width=0.6pt ] (18.75,17.25) rectangle (19,17);
\draw [ color={rgb,255:red,192; green,192; blue,192} , fill={rgb,255:red,192; green,192; blue,192}, line width=0.6pt ] (24.25,7.25) rectangle (24.25,7.25);
\draw [ color={rgb,255:red,192; green,192; blue,192} , fill={rgb,255:red,192; green,192; blue,192}, line width=0.6pt ] (3,21.75) rectangle (3,21.75);

%% ===================================================================
%% STICKSTOFF-PFEILE (N2-Fluss in der Kammer)
%% ===================================================================
\draw [line width=2pt, ->, >=Stealth] (18,10.75) -- (18,13.25) node[pos=0.5, fill=white]{\ce{N2}};
\draw [line width=2pt, ->, >=Stealth] (16.25,13.25) -- (16.25,10.75) node[pos=0.5, fill=white]{\ce{N2}};
\draw [line width=2pt, ->, >=Stealth] (4.5,13.25) -- (4.5,10.75) node[pos=0.5, fill=white]{\ce{N2}};
\draw [line width=2pt, ->, >=Stealth] (6.25,10.75) -- (6.25,13.25) node[pos=0.5, fill=white]{\ce{N2}};

%% ===================================================================
%% PROBENHALTER (schwarzes Rechteck - Sample Matrix)
%% ===================================================================
\draw [ fill={rgb,255:red,0; green,0; blue,0} , line width=1.5pt ] (5.75,8.25) rectangle (16.75,7.25);

%% ===================================================================
%% ADSORPTIONSRÖHRCHEN-SYSTEM (oberer Bereich)
%% ===================================================================
\draw [line width=1.2pt, short] (3.25,19) -- (3.25,17.5);
\draw [ line width=1.2pt ] (3.5,18.25) circle (0.25cm);
\draw [ line width=1.2pt ] (5.25,18.25) circle (0.25cm);
\draw [line width=1.2pt, short] (3.75,19) -- (3.75,17.5);
\draw [line width=1.2pt, short] (5.5,17.5) -- (5.5,19);
\draw [line width=1.2pt, short] (5,19) -- (5,17.5);
\draw [line width=1.2pt, short] (5,19) -- (5.25,19);
\draw [line width=1.2pt, short] (3.75,19) -- (3.5,19);
\draw [line width=1.2pt, short] (4,17.5) -- (4,17);
\draw [line width=1.2pt, short] (4.75,17.5) -- (4.75,17);
\draw [line width=1.2pt, short] (5.25,19) -- (5.25,19.5);
\draw [line width=1.2pt, short] (3.5,19) -- (3.5,19.5);
\draw [line width=1.2pt, short] (5.25,19.5) -- (5.75,19.5);
\draw [line width=1.2pt, short] (3.5,19.5) -- (3,19.5);
\draw [line width=1.2pt, short] (5.75,19.5) -- (5.75,20.75);
\draw [line width=1.2pt, short] (3,19.5) -- (3,20.75);
\draw [line width=1.2pt, short] (5.75,20.75) -- (6.25,20.75);
\draw [line width=1.2pt, short] (6.25,20.5) -- (6,20.5);
\draw [line width=1.2pt, short] (6.25,20.75) -- (6.25,20.5);
\draw [line width=1.2pt, short] (6,20.5) -- (6,19);
\draw [line width=1.2pt, short] (5.5,19) -- (6,19);
\draw [line width=1.2pt, short] (3,20.75) -- (2.5,20.75);
\draw [line width=1.2pt, short] (2.5,20.75) -- (2.5,20.5);
\draw [line width=1.2pt, short] (2.5,20.5) -- (2.75,20.5);
\draw [line width=1.2pt, short] (2.75,20.5) -- (2.75,19);
\draw [line width=1.2pt, short] (2.75,19) -- (3.25,19);
\draw [ line width=1.2pt ] (3.25,19.75) circle (0.25cm);
\draw [ line width=1.2pt ] (5.5,19.75) circle (0.25cm);
\draw [line width=1.2pt, short] (3,20) -- (3.5,20);
\draw [line width=1.2pt, short] (3.5,20) -- (3.5,20.75);
\draw [line width=1.2pt, short] (3.5,20.75) -- (3,20.75);
\draw [line width=1.2pt, short] (5.75,20) -- (5.25,20);
\draw [line width=1.2pt, short] (5.25,20) -- (5.25,20.75);
\draw [line width=1.2pt, short] (5.25,20.75) -- (5.75,20.75);

%% ===================================================================
%% OBERE VERBINDUNGSLEITUNG (vom Adsorptionsröhrchen zur Kammer)
%% ===================================================================
\draw [line width=1.2pt, short] (1.25,17) -- (1.25,18.5);
\draw [line width=1.2pt, short] (1.25,18.5) -- (3.5,18.5);
\draw [line width=1.2pt, short] (5.5,18.5) -- (21.25,18.5);
\draw [line width=1.2pt, short] (21.25,18.5) -- (21.25,18.25);
\draw [line width=1.2pt, short] (3.25,18.5) -- (3.75,18.5);
\draw [line width=1.2pt, short] (3.25,18) -- (3.75,18);
\draw [line width=1.2pt, short] (5,18.5) -- (5.5,18.5);

%% ===================================================================
%% GRAUE FÜLLUNGEN - Verbindungsbereich oben
%% ===================================================================
\draw [ color={rgb,255:red,192; green,192; blue,192} , fill={rgb,255:red,192; green,192; blue,192}, line width=0.2pt ] (1.25,18.5) rectangle (3.25,17);
\draw [ color={rgb,255:red,192; green,192; blue,192} , fill={rgb,255:red,192; green,192; blue,192}, line width=0.2pt ] (2.75,17) rectangle (3.5,16.25);
\draw [ color={rgb,255:red,192; green,192; blue,192} , fill={rgb,255:red,192; green,192; blue,192}, line width=0.2pt ] (3.25,17.5) rectangle (4,17);
\draw [ color={rgb,255:red,192; green,192; blue,192} , fill={rgb,255:red,192; green,192; blue,192}, line width=0.2pt ] (4.75,17.5) rectangle (5.5,17);
\draw [ color={rgb,255:red,192; green,192; blue,192} , fill={rgb,255:red,192; green,192; blue,192}, line width=0.2pt ] (5.5,18.5) rectangle (17.5,17);
\draw [ color={rgb,255:red,192; green,192; blue,192} , fill={rgb,255:red,192; green,192; blue,192}, line width=0.2pt ] (17.5,18.5) rectangle (21.25,18);

%% ===================================================================
%% HAUPTUMRISS (Außenlinien der Kammer - verstärkte Linien 1pt)
%% ===================================================================
\draw [line width=1pt, dashed] (4,17) -- (4.75,17);
\draw [line width=1pt, dashed] (4,17.5) -- (4.75,17.5);
\draw [line width=1pt, short] (4.75,17) -- (17.5,17);
\draw [line width=1pt, short] (5.5,17.5) -- (4.75,17.5);
\draw [line width=1pt, short] (4,17.5) -- (3.25,17.5);
\draw [line width=1pt, short] (2.25,16) -- (2.25,7);
\draw [line width=1pt, short] (2.5,16.75) -- (2.5,7);
\draw [line width=1pt, short] (2.5,7) -- (20,7);
\draw [line width=1pt, short] (2.25,7) -- (2.75,6.25);
\draw [line width=1pt, short] (2.75,6.25) -- (19.75,6.25);
\draw [line width=1pt, short] (19.75,6.25) -- (20.25,7);
\draw [line width=1pt, short] (20.25,7) -- (20.25,16);
\draw [line width=1pt, short] (20.25,16) -- (21.25,16);
\draw [line width=1pt, short] (21.25,16) -- (21.25,16.75);
\draw [line width=1pt, short] (21.25,16.75) -- (20.75,16.75);
\draw [line width=1pt, short] (20,16.25) -- (20,7);
\draw [line width=1pt, short] (20,16.25) -- (19,16.25);
\draw [line width=1pt, short] (19,16.25) -- (19,16.75);
\draw [line width=1pt, short] (19,16.75) -- (18.75,17);
\draw [line width=1pt, short] (18.75,17) -- (18.75,17.25);
\draw [line width=1pt, short] (18.75,17.25) -- (21.25,17.25);
\draw [line width=1pt, short] (21.25,17.25) -- (21.25,17);
\draw [line width=1pt, short] (21.25,17) -- (20.5,17);
\draw [line width=1pt, short] (2.25,16) -- (1.25,16);
\draw [line width=1pt, short] (1.25,16) -- (1.25,16.75);
\draw [line width=1pt, short] (1.25,16.75) -- (1.5,16.75);
\draw [line width=1pt, short] (2,16.75) -- (2.5,16.75);
\draw [line width=1pt, short] (3.5,16.75) -- (3.5,16.25);
\draw [line width=1pt, short] (3.5,16.25) -- (2.75,16.25);
\draw [line width=1pt, short] (2.75,16.25) -- (2.75,17);
\draw [line width=1pt, short] (2.75,17) -- (1.25,17);
\draw [line width=1pt, short] (1.25,17) -- (1.25,18.5);
\draw [line width=1pt, short] (1.25,18.5) -- (3.25,18.5);
\draw [line width=1pt, short] (5.5,18.5) -- (21.25,18.5);
\draw [line width=1pt, short] (21.25,18.5) -- (21.25,18);
\draw [line width=1pt, short] (21.25,18) -- (17.75,18);
\draw [line width=1pt, short] (17.5,17.75) -- (17.5,17);

%% ===================================================================
%% DESORPTIONSKAMMER (grauer Behälter mit gestrichelten Linien)
%% ===================================================================
\draw [ color={rgb,255:red,214; green,214; blue,214} , fill={rgb,255:red,235; green,235; blue,235}, line width=1pt ] (3.5,19) rectangle (5.25,22);
\draw [line width=1pt, short] (3.5,22) -- (3.5,19);
\draw [line width=1pt, short] (3.5,19) -- (5.25,19);
\draw [line width=1pt, short] (5.25,19) -- (5.25,22);
\draw [line width=1pt, dashed] (3.5,19) -- (5.25,22);
\draw [line width=1pt, dashed] (4,19) -- (5.25,21);
\draw [line width=1pt, dashed] (4.5,19) -- (5.25,20.25);
\draw [line width=1pt, dashed] (5,19) -- (5.25,19.5);
\draw [line width=1pt, dashed] (4.75,22) -- (3.5,20);
\draw [line width=1pt, dashed] (3.5,20.75) -- (4.25,22);
\draw [line width=1pt, dashed] (3.5,21.5) -- (3.75,22);

\draw [line width=1pt, short] (3.25,18.5) -- (3.25,17.5);
\draw [line width=1pt, short] (5.5,18.5) -- (5.5,17.5);
\draw [line width=1pt, short] (4.75,17.5) -- (4.75,17);
\draw [line width=1pt, short] (4,17.5) -- (4,17);

%% ===================================================================
%% BESCHRIFTUNGEN UND PFEILE
%% ===================================================================
% Probenmatrix-Pfeil (von unten nach oben zur Kammer)
\draw [line width=1.5pt, ->, >=Stealth] (4.5,15.75) -- (4.5,18.5)node[pos=0.05,below, xshift=12, fill=white]{\Huge Probenmatrix };

% N2-Ausgang oben (aus Desorptionskammer)
\draw [line width=1pt, ->, >=Stealth] (4.5,22.25) -- (4.5,24)node[pos=0.95,above, fill=white]{\Huge \ce{N2}};

\draw [line width=1pt, dashed] (3.5,22) -- (5.25,22);

%% ===================================================================
%% HEIZBLOCK (unterer Bereich)
%% ===================================================================
\draw [ color={rgb,255:red,192; green,192; blue,192} , fill={rgb,255:red,192; green,192; blue,192}, line width=1pt ] (2,6) rectangle (20.5,4.5);
\draw [line width=1pt, short] (2,6) -- (20.5,6);
\node [font=\Huge] at (11.25,5.25) {Heizblock};

%% ===================================================================
%% EXTERNES FLIESSSCHEMA (Ventil und Durchflussmesser rechts)
%% ===================================================================
\draw [line width=1pt, ->, >=Stealth] (25.3,17.65) -- (18.25,17.65);
\draw (27.7,17.65) -- (30.2,17.65)node[pos=-0.475, below, yshift=-20, fill=white]{\Huge V1};
\draw (30.2,17.65) -- (30.2,11);
\pic[scale=5.5] (valve) at (26.5,17.65) {valve=main};

% Durchflussmesser-Symbol
\draw (34,11) circle (1.8cm);
\draw (30.2,11) -- (32.8,11);
\draw (35.2,11) -- (38.5,11)node[pos=1, fill=white]{\Huge \ac{dvdt}$_{\ce{N2}}$};%$\dot{V}_{\ce{N2}}$};
\draw (32.8,11) -- (33.5,12);
\draw (35.2,11) -- (34.5,10); %node[pos=2.2,below, xshift=9, fill=white]{\Huge H1};
\draw (33.5,12) -- (34.5,10);
\draw [line width=2.5, ->, >=Stealth] (36.5,14) -- (31.5,8);

%% ===================================================================
%% WÄRME-/DAMPFPFEILE (gestrichelt, vom Probenhalter aufsteigend)
%% ===================================================================
\draw [line width=1pt, ->, >=Stealth, dashed] (6.25,8.5) -- (5.25,9.5);
\draw [line width=1pt, ->, >=Stealth, dashed] (7.25,8.5) -- (7,9.75);
\draw [line width=1pt, ->, >=Stealth, dashed] (8.75,8.75) -- (8.75,9.75);
\draw [line width=1pt, ->, >=Stealth, dashed] (16.25,8.5) -- (17.25,9.5);
\draw [line width=1pt, ->, >=Stealth, dashed] (15.25,8.5) -- (15.75,9.75);
\draw [line width=1pt, ->, >=Stealth, dashed] (13.75,8.75) -- (13.75,9.75);

% Prüfkörper-Beschriftung
\draw [line width=1pt, ->, >=Stealth] (11.25,10.75) -- (10.5,8.5)node[pos=0,above, fill=white]{\Huge Prüfkörper};

%% ===================================================================
%% ZUSÄTZLICHE DETAIL-FÜLLUNGEN
%% ===================================================================
\draw [ color={rgb,255:red,192; green,192; blue,192} , fill={rgb,255:red,192; green,192; blue,192}, line width=1pt ] (3.25,16.5) -- (3,16.5) -- (3.5,17) -- (3.75,17) -- cycle;
\draw [ color={rgb,255:red,192; green,192; blue,192} , fill={rgb,255:red,192; green,192; blue,192}, line width=1pt ] (17.25,17.5) -- (17,17.5) -- (17.5,18) -- (17.75,18) -- cycle;
\draw [line width=1pt, short] (3.5,16.75) -- (3.75,17);
\draw [line width=1pt, short] (3.75,17) -- (4,17);
\draw [line width=1pt, short] (17.5,17.75) -- (17.75,18);

%% ===================================================================
%% WEITERE BESCHRIFTUNGEN
%% ===================================================================
% Kammerverschluss
\draw [line width=1pt, ->, >=Stealth] (15,20.75) -- (11.5,18.75)node[pos=0,above, fill=white]{\Huge Kammerverschluss};

\draw [line width=1pt, dashed] (17.5,17) -- (18.75,17);
\draw [line width=1pt, dashed] (21.25,18) -- (21.25,17.25);

% Adsorptionsröhrchen
\draw [line width=1pt, ->, >=Stealth] (0,22.5) -- (3.25,21.5)node[pos=0,above, fill=white]{\Huge Adsorbtionsröhrchen};

\draw [ fill={rgb,255:red,192; green,192; blue,192} , line width=1pt ] (21.5,22.75) circle (0cm);

%% ===================================================================
%% SEITLICHE STÜTZSTRUKTUR (links und rechts)
%% ===================================================================
% Linke Seite
\draw [ color={rgb,255:red,192; green,192; blue,192} , fill={rgb,255:red,192; green,192; blue,192}, line width=1pt ] (1.5,4.75) circle (0.25cm);
\draw [ color={rgb,255:red,192; green,192; blue,192} , fill={rgb,255:red,192; green,192; blue,192}, line width=1pt ] (1.5,5.75) rectangle (2,4.5);
\draw [ color={rgb,255:red,192; green,192; blue,192} , fill={rgb,255:red,192; green,192; blue,192}, line width=1pt ] (1.25,5.75) rectangle (1.5,4.75);
\draw [ color={rgb,255:red,192; green,192; blue,192} , fill={rgb,255:red,192; green,192; blue,192}, line width=1pt ] (1.25,15.75) rectangle (2,5.5);
\draw [line width=1pt, short] (1.25,15.75) -- (1.25,4.75);
\draw [line width=1pt, short] (1.25,15.75) -- (2,15.75);
\draw [line width=1pt, short] (2,15.75) -- (2,6);
\draw [line width=1pt, short] (1.25,4.75) .. controls (1.25,4.5) and (1.25,4.5) .. (1.5,4.5);

% Rechte Seite
\draw [ color={rgb,255:red,192; green,192; blue,192} , fill={rgb,255:red,192; green,192; blue,192}, line width=1pt ] (21,4.75) circle (0.25cm);
\draw [ color={rgb,255:red,192; green,192; blue,192} , fill={rgb,255:red,192; green,192; blue,192}, line width=1pt ] (20.5,6) rectangle (21,4.5);
\draw [ color={rgb,255:red,192; green,192; blue,192} , fill={rgb,255:red,192; green,192; blue,192}, line width=1pt ] (21,6) rectangle (21.25,5);
\draw [ color={rgb,255:red,192; green,192; blue,192} , fill={rgb,255:red,192; green,192; blue,192}, line width=1pt ] (21,5.25) rectangle (21.25,4.75);
\draw [ color={rgb,255:red,192; green,192; blue,192} , fill={rgb,255:red,192; green,192; blue,192}, line width=1pt ] (21.25,6) rectangle (20.5,15.75);
\draw [line width=1pt, short] (1.5,4.5) -- (21,4.5);
\draw [line width=1pt, short] (20.5,6) -- (20.5,15.75);
\draw [line width=1pt, short] (20.5,15.75) -- (21.25,15.75);
\draw [line width=1pt, short] (21.25,15.75) -- (21.25,4.75);
\draw [line width=1pt, short] (21.25,4.75) .. controls (21.25,4.5) and (21.25,4.5) .. (21,4.5);

\end{circuitikz}
}%
\caption{Schematische Darstellung einer Micro-Chamber. Angepasst aus Lit. \cite{micro_chamber_figure}}
\label{fig:my_label}
\end{figure}